\section{Analysis}
\label{sec:analysis}

We examine four dimensions of the preliminary \IIRC results from
Table~\ref{tab:iirc-main}, supplemented by the \TwoWikiShort calibration
surface and targeted ablations.

\paragraph{Precision--recall tradeoff.}
The three operating points in Table~\ref{tab:iirc-main} expose a
clear precision--recall separation.
The controller achieves $\mathrm{P} = 0.533$ and $\mathrm{R} = 0.475$
with a compact selection of 1.75~nodes per case.
Dense retrieval at a matched node count (tokens-512, also 1.75~nodes)
reaches $\mathrm{P} = 0.458$ and $\mathrm{R} = 0.389$---lower on
both axes.
Both selectors operate in the same region of the node-count spectrum,
so the controller's advantage is not an artifact of retrieving more
documents; it selects \emph{better} documents by following link-context
signals to bridge evidence that embedding similarity alone misses.

Under ratio-controlled budgets, dense retrieval occupies a qualitatively
different region: $\mathrm{P} = 0.022$ and $\mathrm{R} = 0.622$ with
64~nodes per case.
The recall gain over fixed-token dense ($+0.233$) is real but purchased
at a $20{\times}$ precision collapse, yielding $\Fempty
= 0.042$.
This configuration retrieves roughly one in every thousand corpus
documents, yet fewer than one in fifty of the retrieved documents is
a gold support page.
The \TwoWikiShort calibration surface corroborates this pattern at a smaller
scale: the single-path walk gains recall ($+0.128$) over dense at tokens-256
while matching its precision ($0.310$ vs.\ $0.311$)
(Table~\ref{tab:2wiki-calibration}), demonstrating that link-context expansion
does not merely trade one metric for the other.

\paragraph{Budget sensitivity.}
A striking feature of the preliminary surface is the controller's
\emph{budget independence}: its $\Fempty$, precision,
recall, and node count are identical across ratio-0.01 through ratio-1.0.
This occurs because the controller's walk decisions---which links to
follow, when to stop---are functions of query--link-context relevance,
not of the available budget.
The ratio only governs the downstream budget-backfill stage, and on these
cases the controller's walk already exhausts the useful link-context
neighborhood before backfill is invoked.
In effect, the controller is \textbf{self-limiting}: it selects what the
query needs and stops, regardless of how much budget remains.

Dense retrieval under ratio budgets exhibits the opposite behavior.
As the budget ratio grows from 0.01 to~1.0, the dense selector fills its
allocation with progressively more documents---up to 64~per case at
$\rho = 0.01$ on full-\IIRC---because its only stopping criterion is
budget exhaustion.
This \textbf{budget-filling} pattern explains the precision collapse
observed in Section~\ref{sec:results:iirc-ratio}: more budget means more noise,
not more signal.
The asymmetry has practical implications: in deployment, a self-limiting
selector degrades gracefully when given an overly generous budget, while
a budget-filling selector requires careful budget tuning to avoid
drowning downstream components in irrelevant passages.

\paragraph{When does the controller help versus hurt?}
On the 20-case \IIRC pilot, the controller ($\Fempty =
0.471$) outperforms dense at tokens-512 ($0.408$) by $+0.063$;
projected 100-case estimates narrow this gap somewhat
(\selcontroller $\approx 0.44$, \seldense $0.337$) but preserve
the direction.
This gain requires qualification along two axes: walk strategy and
model capability.

First, the controller's quality-aware stopping is what separates it
from the walk-only selectors.
Non-controller walks using the best non-LLM scorer
(\selwalk with sentence-transformer, lookahead-2,
\stfuture) achieved only modest improvements over dense
on \TwoWikiShort ($+0.028$ F1, Table~\ref{tab:2wiki-calibration}).
On earlier 100-case \IIRC runs under fixed-token budgets, non-LLM
single-path walks \emph{underperformed} dense
($\Fempty = 0.273$ vs.\ $0.319$ at tokens-256).
The non-LLM scorers can rank links by shallow overlap or embedding
similarity, but they cannot reason about whether a candidate document
would bridge to unseen evidence---precisely the capability the LLM
controller provides.

Second, controller quality is model-capability-dependent.
\controllerlm achieves $\Fempty \approx 0.44$ on projected 100-case
estimates, while replacing it with \controllersmall as the controller
LLM yields $\Fempty = 0.235$---far below dense and
below the non-LLM walk.
% TODO: confirm \controllerlm F1 on 100-case canonical surface.
Smaller models lack the reasoning capacity to evaluate link context
against the query reliably: they follow links that surface-match the
query terms without assessing whether the target document would supply
missing evidence.
This negative result bounds the method's applicability:
controller-guided selection requires a sufficiently capable LLM,
and the cost--quality tradeoff (Section~\ref{sec:results:iirc-fixed}) must be
evaluated against the target deployment budget.

\paragraph{Selectivity analysis.}
The controller selects a mean of 1.75~nodes per case, compared to the
gold oracle's 2.55~nodes.
It is slightly undershooting the gold set size---selecting 69\% of the
oracle's node count---which accounts for part of the recall gap
($\mathrm{R} = 0.475$ vs.\ $1.0$).
However, the more informative comparison is with dense under ratio
budgets, which selects 64~nodes per case: $36{\times}$ more than the
controller for dramatically worse F1 ($0.042$ vs.\ $0.471$).
The controller achieves 47.1\% of the oracle's $\Fempty$
with 69\% of the oracle's node count, while dense at ratio-0.01 achieves
4.2\% of the oracle's F1 with $25{\times}$ the oracle's node count.
% TODO: compute corpus-mass compression ratio (selected tokens / total tokens)
% for controller vs dense on 100-case surface.

This selectivity gap reflects a difference in how the two selectors
cope with uncertainty.
The dense selector hedges by including every document above a similarity
threshold, producing a large candidate set that is likely to contain the
gold documents but buries them in noise.
The controller commits to a small number of high-confidence selections by
reasoning about link context, accepting the risk of undershooting in
exchange for a dramatically higher signal-to-noise ratio.
For downstream consumers---whether a reader model, a GraphRAG
pipeline, or a human analyst---the controller's compact, high-precision
selection is more useful per token than the dense selector's
exhaustive retrieval, even when the latter achieves higher recall.

\paragraph{Walk structure ablation.}
Table~\ref{tab:walk-structure} compares alternative graph-search strategies
on \TwoWikiShort (30 cases, budget 256).
Each strategy operates in its natural regime: beam and A* use 3~seeds and
3~hops to exploit their broader frontier, while the single-path walk uses
1~seed and 2~hops.
Broader graph search (beam, A*) achieves high recall---matching or
exceeding the single-path walk's recall---but collapses precision,
producing substantially lower $\Fempty$.
The single-path walk's advantage comes from its precision discipline:
it adds only one document per step and stops when scores drop, while
beam and A* expand aggressively and dilute the selection with
marginally relevant documents.
Even the dense baseline outperforms beam and A* in F1, confirming that
unconstrained graph exploration is counterproductive under tight budgets.
These results motivate the design choice of constrained, quality-aware
walking over unconstrained graph exploration.

\begin{table}[t]
\centering
\caption{Walk structure ablation on \TwoWikiShort (30 cases, token budget = 256).
Beam and A* use top-3 seeds and 3 hops; single-path uses top-1 seed and 2 hops.}
\label{tab:walk-structure}
\small
\setlength{\tabcolsep}{4pt}
\begin{tabular}{@{}lccccc@{}}
\toprule
\textbf{Walk Strategy} & \textbf{Seeds} & \textbf{Hops}
  & \textbf{$\Fempty$} & \textbf{Prec.} & \textbf{Recall} \\
\midrule
Single-path (ST+LLM) & 1 & 2 & \textbf{.472} & \textbf{.378} & \textbf{.750} \\
Beam (ST+LLM)        & 3 & 3 & .311 & .203 & .733 \\
Beam (overlap)        & 3 & 3 & .227 & .156 & .458 \\
A* (ST+LLM)          & 3 & 3 & .191 & .112 & .733 \\
A* (overlap)          & 3 & 3 & .147 & .092 & .483 \\
\midrule
Dense (ST, fill)      & 1 & 0 & .412 & .334 & .617 \\
\bottomrule
\end{tabular}
\end{table}

\paragraph{Walk vs.\ fill decomposition.}
Budget fill---using remaining token headroom for additional dense
retrieval after the walk completes---contributes meaningfully to all
selectors on \TwoWikiShort (30 cases, budget 256).
Without fill, dense retrieval leaves 26.7\% of selections empty,
and $\Fempty$ drops from $0.354$ to $0.302$; fill eliminates
empty selections and raises recall from $0.217$ to $0.408$, though
precision falls from $0.727$ to $0.331$ as the additional documents
are lower-confidence.
For the LLM-guided walk, fill raises $\Fempty$ from $0.317$ to
$0.367$ by backfilling remaining budget headroom with dense retrieval
($\mathrm{R}$: $0.317 \to 0.475$; $\mathrm{P}$: $0.385 \to 0.324$).
However, the walk's contribution is not subsumed by fill:
the walk discovers bridge documents that dense retrieval cannot reach
via embedding similarity alone, while fill provides coverage insurance
for cases where the link graph offers no productive path.

\paragraph{Edge-scorer ablation.}
Table~\ref{tab:edge-scorer-ablation} reports the edge-scorer ablation
on \TwoWikiShort (100 cases, budget 128, single-path walk).
The sentence-transformer scorer with \stfuture weighting and
lookahead~2 achieves the highest $\Fempty = 0.394$, outperforming
both the balanced profile ($0.385$) and direct-heavy profile ($0.356$)
at the same lookahead depth.
Reducing lookahead from 2 to~1 drops balanced-profile F1 from $0.385$
to $0.354$, confirming that two-step lookahead provides meaningful
signal about the downstream value of a link.
The gap between the best ST scorer ($0.394$) and both overlap baselines
($\leq 0.352$) demonstrates that embedding-based scoring provides
meaningful signal beyond lexical matching, though the absolute margin
is modest ($+0.042$).
This modest gap explains why the LLM controller---which can reason
about link context rather than merely scoring it---delivers a
substantially larger gain over the non-LLM walk.
